\section{Sparse coordinate strips}
\label{sec:area_law_sparse_strips}

\begin{definition}[Coordinate residue strips]
\label{def:area_law_residue_strips}
Let $F\subset\mathbb Z^2$ be finite and let $m\ge1$ be an integer.
For residues $a,b\in\{0,\ldots,m-1\}$, define
\begin{align}
 B_{a,b}(F)&=\{z\in F:z_1\equiv a\pmod m\ \text{or}\ z_2\equiv b\pmod m\}.
 \label{eq:area_law_residue_strips}
\end{align}
\end{definition}

\begin{theorem}[Sparse choice of residue strips]
\label{thm:area_law_sparse_belt_shift}
\lean{TNLean.PEPS.AreaLaw.Geometry.exists_sparse_belt_shift}
\leanok
\uses{def:area_law_residue_strips}
For every finite $F\subset\mathbb Z^2$ and integer $m\ge1$, there are
residues $a,b\in\{0,\ldots,m-1\}$ such that
\begin{align}
 m|B_{a,b}(F)|&\le2|F|.
 \label{eq:area_law_sparse_belt_shift}
\end{align}
This is the finite averaging step in the sparse-strip construction of
\cite[Section~11]{OpenAI2026AreaLaw}.
% Source revision: openai/math@adc7f1241b42e322a6451854ab7e4b4c146bf78a,
% build/sections/10-geometry.tex, lines 210--237, geometry:belt-count.
% This auxiliary theorem does not establish the layer-cell count or the
% final dyadic decay estimate in geometry:belt-count.
\end{theorem}
\begin{proof}\leanok
For any map $f:F\to\{0,\ldots,m-1\}$, its $m$ fibres partition $F$.
One fibre therefore has at most $\lfloor|F|/m\rfloor$ elements.
Apply this assertion to the two coordinate residue maps and choose their
residues independently. Each selected strip contains at most
$\lfloor|F|/m\rfloor$ cells, so their union satisfies
$|B_{a,b}(F)|\le2\lfloor|F|/m\rfloor$. Multiplying by $m$ and using
$m\lfloor|F|/m\rfloor\le|F|$ gives
(\ref{eq:area_law_sparse_belt_shift}). The argument includes $F=\varnothing$
and $m=1$.
\end{proof}

\begin{definition}[Boundary endpoints and dyadic children]
\label{def:area_law_cell_counts}
\lean{
  TNLean.PEPS.AreaLaw.Geometry.boundaryEndpoints,
  TNLean.PEPS.AreaLaw.Geometry.dyadicChildren
}
\leanok
\uses{def:area_law_hamiltonian}
Let $Z$ be the set of integer endpoints of the crossing edges of a cut.
For a finite set $P\subset\mathbb Z^2$ of parent indices, its dyadic child
indices are
\begin{align}
 \operatorname{Ch}(P)&=
 \{2z+\varepsilon:z\in P,\ \varepsilon\in\{0,1\}^2\}.
 \label{eq:area_law_dyadic_children}
\end{align}
These are the four children of each parent in a translated dyadic mesh.
\end{definition}

\begin{theorem}[Finite cell counts]
\label{thm:area_law_cell_counts}
\lean{
  TNLean.PEPS.AreaLaw.Geometry.card_ambientDilation_le,
  TNLean.PEPS.AreaLaw.Geometry.card_boundaryEndpoints_le,
  TNLean.PEPS.AreaLaw.Geometry.card_dyadicChildren
}
\leanok
\uses{def:area_law_cell_counts,def:area_law_template}
For every finite $T\subset\mathbb Z^2$, integer $r\ge0$, finite parent
set $P$, and finite induced domain with cut $A$, one has
\begin{align}
 |(T+[-r,r]^2)\cap\mathbb Z^2|&\le(2r+1)^2|T|,
 \label{eq:area_law_dilation_count}\\
 |Z|&\le2|\partial_\Lambda A|,
 \label{eq:area_law_endpoint_count}\\
 |\operatorname{Ch}(P)|&=4|P|.
 \label{eq:area_law_children_count}
\end{align}
These cardinality estimates are ingredients in the layer count of
\cite[Section~11]{OpenAI2026AreaLaw}; no connectedness assumption is needed.
\end{theorem}
\begin{proof}\leanok
Each point of $T$ contributes an integer square with $(2r+1)^2$ sites;
counting their union gives the first bound. Each crossing edge contributes
at most two endpoints, giving the second. In each coordinate the remainder
modulo two identifies the child label, and division by two then identifies
the parent index. Thus distinct parent and child-label pairs give distinct
indices, proving the last equality. All three assertions include empty sets.
\end{proof}
